%=================================================================================
\section[MD-LDP]{Multi-Distance Lidar Data Processing}
%=================================================================================
\begin{frame}{The idea in one picture}
	\centering
	\includegraphics[width=\linewidth]{Figures/MD_LDP_concept}
	\vfill
	\begin{block}<1->{Same interface as \texttt{LDP\_v3}}
		One call per time step, only the lidar signals (line-of-sight speed, valid flag, beam ID) of the current
		and past time steps, but \emph{all} gates (multi-distance idea \cite{Schlipf2015e}).
		Two outputs from the same buffer: one for each discipline.
	\end{block}
\end{frame}
%---------------------------------------------------------------------------------
\begin{frame}{Key idea 1: blade ghost avoidance}
	\begin{columns}[T]
		\column{7.0cm}
		\begin{block}<1->{What goes wrong in the baseline}
			\begin{itemize}
				\item mean over the beams that happen to be valid
				\item when a blade blocks a beam, the \emph{composition} of the mean changes
			\end{itemize}
		\end{block}
		\column{7.0cm}
		\begin{block}<1->{What we do}
			\begin{itemize}
				\item each beam and gate keeps its own last valid information
				\item a gate is always represented by all 4 beams
			\end{itemize}
		\end{block}
	\end{columns}
	\begin{alertblock}<1->{Effect}
		With the baseline settings only this step: SDES $1.738\,D \rightarrow 1.657\,D$
	\end{alertblock}
\end{frame}
%---------------------------------------------------------------------------------
\begin{frame}{Key idea 2: think in rotor time}
	\begin{columns}[T]
		\column{7.0cm}
		\begin{block}<1->{Every measurement gets its arrival time}
			\begin{itemize}
				\item the air measured at a gate reaches the rotor seconds later
				\item each measurement is stored with: \emph{when} it will arrive at the rotor
				\item the output is read from this buffer at the desired rotor time
			\end{itemize}
		\end{block}
		\column{7.0cm}
		\begin{block}<1->{The preview pays twice}
			\begin{itemize}
				\item the \emph{next} measurement of a beam is usually taken before the current one arrives
				\item  interpolation instead of holding
				\item all gates aligned to the same rotor time can be combined
			\end{itemize}
		\end{block}
	\end{columns}
	\begin{alertblock}<1->{Largest single gain}
		Combining the aligned distances: SDES $1.590\,D \rightarrow 1.270\,D$ (more than all other steps together)
	\end{alertblock}
\end{frame}
%---------------------------------------------------------------------------------
\begin{frame}{Key idea 3: which speed transports the turbulence?}
	\begin{columns}[T]
		\column{7.0cm}
		\begin{block}<1->{Space instead of time}
			\begin{itemize}
				\item a long time average at a fixed point mostly contains air that has \emph{already passed}
				\item the measured volume upstream contains the air that is \emph{about to arrive}
			\end{itemize}
		\end{block}
		\column{7.0cm}
		\begin{block}<1->{Physics on top}
			\begin{itemize}
				\item induction zone in front of the rotor
				\item height profile of the wind
				\item small effects on the SDES, but consistent and cheap
			\end{itemize}
		\end{block}
	\end{columns}
	\begin{alertblock}<1->{Effect}
		Transport speed from space instead of a time filter: SDES $1.222\,D \rightarrow 1.157\,D$
	\end{alertblock}
\end{frame}
%---------------------------------------------------------------------------------
\begin{frame}{Key idea 4: one buffer, two outputs}
	\begin{columns}[T]
		\column{7.0cm}
		\begin{block}<1->{Discipline 1: preview quality}
			\begin{itemize}
				\item read at the current rotor time
				\item no low-pass filter
				\item the SDES rewards coherent content
			\end{itemize}
		\end{block}
		\column{7.0cm}
		\begin{block}<1->{Discipline 2: feedforward control}
			\begin{itemize}
				\item read a few seconds \emph{ahead} of the rotor
				\item adaptive low-pass filter: only what the turbine can follow
				\item the preview compensates the delay of the filter
			\end{itemize}
		\end{block}
	\end{columns}
	\begin{alertblock}<1->{Lesson learned}
		A filtered signal that arrives too late is worse than no lidar at all:
		late feedforward fights the feedback controller and increases the pitch activity.
	\end{alertblock}
\end{frame}
